\par\noindent\textbf{(1)} 把粒子相互作用写成 \(\sum_{i<j}u(r_{ij})\)，经典构型权重为 \(e^{-\beta\sum u}\)。定义Mayer函数 \(f(r)=e^{-\beta u(r)}-1\)。二阶簇积分由两粒子相对位移积分给出 \(b_2=\frac12\int f(\boldsymbol r)\,d^3r\)，与理想气体的 \(N!\) 计数区分开。
\par\noindent\textbf{(2)} 簇展开消去逸度后得到 \(B_2=-b_2=-\frac12\int_{\mathbb R^3}[e^{-\beta u(r)}-1]\,d^3r\)。对各向同性势化成 \(B_2=-2\pi\int_0^\infty[e^{-\beta u(r)}-1]r^2dr\)。若吸引尾部不使积分收敛，普通低密Virial级数不能直接使用。
\par\noindent\textbf{(3)} 硬球模型 \(u(r)=+\infty\) 对 \(r<d\)、其余为零，故 \(f=-1\) 在排除球内、其余为零。积分给 \(B_2=2\pi\int_0^d r^2dr=2\pi d^3/3\)，是单粒子硬球体积 \(\pi d^3/6\) 的四倍。正号表示排除体积使压强高于理想气体。
\par\noindent\textbf{(4)} 若某温区存在明显吸引势 \(u<0\)，则相应区间 \(e^{-\beta u}-1>0\)，对 \(B_2\) 给负贡献；在Boyle温度两类贡献可能抵消。必须要求 \(|B_2|n\ll1\) 且更高阶簇积分可以忽略，不能用第二Virial近似描述凝聚临界区。
